\documentclass{article}
\usepackage[UTF8]{ctex}
\usepackage{amsmath,amssymb,amsfonts}
\usepackage{enumitem}
\usepackage{geometry}
\geometry{a4paper, margin=1in}
% === 错题本排版优化（不参与题目标识） ===
\setlength{\parindent}{0pt}
\setlength{\parskip}{0.35em plus 0.10em minus 0.08em}
\setlength{\emergencystretch}{3em}
\setlist{itemsep=0.25em, topsep=0.40em, parsep=0pt}
\allowdisplaybreaks[2]
% === 错题本排版优化结束 ===
\begin{document}
\section*{东南大学2025年高等代数真题}
\section*{80．东南大学2025年高等代数试卷（回忆版）}
\begin{itemize}

\subsection{题目1}
  \item[1．] 设 \(A=\left(a_{i j}\right)_{3 \times 3}\) 的各元素的代数余子式之和为 1 ，求
\[
\left|\begin{array}{lll}
1+a_{11} & 1+a_{12} & 1+a_{13} \\
1+a_{21} & 1+a_{22} & 1+a_{23} \\
1+a_{31} & 1+a_{32} & 1+a_{33}
\end{array}\right|
\]

\subsection{题目2}
  \item[2．] 多项式 \(f(x)\) 无重因式，\(f^{\prime}(x)\) 是 \(f(x)\) 的微商，\(A\) 是 3 阶方阵，求
\[
r\binom{f(A)}{f^{\prime}(A)} .
\]
\end{itemize}

\subsection{题目3}
3．设

\[
A=\left(\begin{array}{lll}
a & * & * \\
* & 3 & * \\
* & * & 2
\end{array}\right), B=\left(\begin{array}{lll}
* & * & * \\
* & * & c \\
* & * & *
\end{array}\right) \text {. }
\]

（1）求 \(A\) 的若尔当标准形与不变因子．\\
（2）\(A\) 与 \(B\) 何时相似？

\subsection{题目4}
4．\(n\) 维欧氏空间 \(V\) 上的变换 \(f\) 满足对任意的 \(\alpha \in V\) ，有

\[
f(\alpha)=a \alpha-b(\alpha, \eta) \eta,\|\eta\|=\sqrt{3} .
\]

\(a, b\) 取何值时，\(f\) 是正交变换？

\subsection{题目5}
5．\(n\) 阶方阵 \(A, B\) 满足 \((A+B)^{-1}=A^{-1}+B^{-1}\) 。

\begin{itemize}
  \item[（1）] 证明：\(|A+B|^{2}=|A||B|\) ．
  \item[（2）] \(A, B\) 是实矩阵，证明：\(|A|=|B|\) 。
  \item[（3）] \(A, B\) 是复矩阵，\(|A|=|B|\) 是否成立？
\end{itemize}

\subsection{题目6}
6．\(n\) 阶方阵 \(A\) 满足 \(A_{11} \neq 0, b\) 是 \(n\) 维非零向量，证明：\(A X=0\) 有无穷解的充要条件是 \(A^{*} X=b\) 有解．

\subsection{题目7}
7．\(A\) 是 \(m \times n\) 矩阵。

\begin{itemize}
  \item[（1）] 存在矩阵 \(B\) 满足 \(A B A=A, B A B=B\) 。
  \item[（2）] 若 \(B\) 满足 \(A B A=A\) 且 \(B A B=B, K(A)=\left\{A X=0 \mid X \in P^{n}\right\}, R(B)=\left\{B Y \mid Y \in P^{m}\right\}\) ，证明：
\[
P^{n}=K(A) \oplus R(B) .
\]
\end{itemize}

\subsection{题目8}
8．\(A, B\) 是 \(n\) 阶方阵。

\begin{itemize}
  \item[（1）] \(\operatorname{tr}(A B)\) 是否等于 \(\operatorname{tr}(B A)\) ？
  \item[（2）] \(A, B, A-B\) 正定，证明： \(\operatorname{tr}\left(A^{2}-B^{2}\right)>0\) ．
\end{itemize}

\subsection{题目9}
9．一个关于什么矩阵存在性的问题．
\end{document}
